%% Generated by Sphinx.
\def\sphinxdocclass{report}
\documentclass[letterpaper,10pt,polish]{sphinxmanual}
\ifdefined\pdfpxdimen
   \let\sphinxpxdimen\pdfpxdimen\else\newdimen\sphinxpxdimen
\fi \sphinxpxdimen=.75bp\relax
\ifdefined\pdfimageresolution
    \pdfimageresolution= \numexpr \dimexpr1in\relax/\sphinxpxdimen\relax
\fi
\newdimen\sphinxremdimen\sphinxremdimen = 10pt
%% let collapsible pdf bookmarks panel have high depth per default
\PassOptionsToPackage{bookmarksdepth=5}{hyperref}

\PassOptionsToPackage{booktabs}{sphinx}
\PassOptionsToPackage{colorrows}{sphinx}

\PassOptionsToPackage{warn}{textcomp}
\usepackage[utf8]{inputenc}
\ifdefined\DeclareUnicodeCharacter
% support both utf8 and utf8x syntaxes
  \ifdefined\DeclareUnicodeCharacterAsOptional
    \def\sphinxDUC#1{\DeclareUnicodeCharacter{"#1}}
  \else
    \let\sphinxDUC\DeclareUnicodeCharacter
  \fi
  \sphinxDUC{00A0}{\nobreakspace}
  \sphinxDUC{2500}{\sphinxunichar{2500}}
  \sphinxDUC{2502}{\sphinxunichar{2502}}
  \sphinxDUC{2514}{\sphinxunichar{2514}}
  \sphinxDUC{251C}{\sphinxunichar{251C}}
  \sphinxDUC{2572}{\textbackslash}
\fi
\usepackage{cmap}
\usepackage[T1]{fontenc}
\usepackage{amsmath,amssymb,amstext}
\usepackage{babel}



\usepackage{tgtermes}
\usepackage{tgheros}
\renewcommand{\ttdefault}{txtt}



\usepackage[Sonny]{fncychap}
\ChNameVar{\Large\normalfont\sffamily}
\ChTitleVar{\Large\normalfont\sffamily}
\usepackage{sphinx}

\fvset{fontsize=auto}
\usepackage{geometry}


% Include hyperref last.
\usepackage{hyperref}
% Fix anchor placement for figures with captions.
\usepackage{hypcap}% it must be loaded after hyperref.
% Set up styles of URL: it should be placed after hyperref.
\urlstyle{same}

\addto\captionspolish{\renewcommand{\contentsname}{Contents:}}

\usepackage{sphinxmessages}
\setcounter{tocdepth}{1}



\title{Partycjonowanie Danych}
\date{27 cze 2026}
\release{}
\author{Dominiak Bystrzak}
\newcommand{\sphinxlogo}{\vbox{}}
\renewcommand{\releasename}{}
\makeindex
\begin{document}

\ifdefined\shorthandoff
  \ifnum\catcode`\=\string=\active\shorthandoff{=}\fi
  \ifnum\catcode`\"=\active\shorthandoff{"}\fi
\fi

\pagestyle{empty}
\sphinxmaketitle
\pagestyle{plain}
\sphinxtableofcontents
\pagestyle{normal}
\phantomsection\label{\detokenize{index::doc}}


\sphinxAtStartPar
Add your content using \sphinxcode{\sphinxupquote{reStructuredText}} syntax. See the
\sphinxhref{https://www.sphinx-doc.org/en/master/usage/restructuredtext/index.html}{reStructuredText}
documentation for details.

\sphinxstepscope

\sphinxAtStartPar
/Author: {[}Damian Dominiak, Michal Bystrzak{]} / Date: 2026\sphinxhyphen{}05\sphinxhyphen{}15 / Version: 1.0 /


\chapter{Rozdział 1. Wprowadzenie}
\label{\detokenize{rozdzial_1/index:rozdzial-1-wprowadzenie}}\label{\detokenize{rozdzial_1/index::doc}}

\section{1.1. Repozytoria wykorzystane w projekcie}
\label{\detokenize{rozdzial_1/index:repozytoria-wykorzystane-w-projekcie}}
\sphinxAtStartPar
Poniżej znajdują się odnośniki do repozytoriów zawierających kody źródłowe oraz dokumentację sprawozdania:
\begin{itemize}
\item {} 
\sphinxAtStartPar
\sphinxstylestrong{Sprawozdanie główne i dokumentacja (Sphinx):} \sphinxhref{https://github.com/...}{{[}Wklej link do repozytorium GitHub{]}}

\item {} 
\sphinxAtStartPar
\sphinxstylestrong{Badania literaturowe (Michał Bystrzak):} \sphinxhref{https://github.com/...}{{[}Wklej link do podłączonego repozytorium{]}}

\item {} 
\sphinxAtStartPar
\sphinxstylestrong{Eksperymenty i skrypty bazodanowe:} \sphinxhref{https://github.com/...}{{[}Wklej link do repozytorium ze skryptami .py{]}}

\end{itemize}


\section{1.2. Opis tematyki projektu}
\label{\detokenize{rozdzial_1/index:opis-tematyki-projektu}}
\sphinxAtStartPar
Przedmiotem niniejszego projektu jest zaprojektowanie, zdefiniowanie oraz obsługa relacyjnej bazy danych dla systemu zarządzania serwisem sprzętu IT. Aplikacja docelowa ma za zadanie obsługiwać pełen cykl życia zlecenia serwisowego \textendash{} od momentu przyjęcia sprzętu od klienta, poprzez diagnozę i zużycie części magazynowych, aż po ostateczne rozliczenie i wydanie naprawionego urządzenia.

\sphinxAtStartPar
W ramach zajęć laboratoryjnych przygotowano schematy (koncepcyjny, logiczny i fizyczny), przeprowadzono normalizację bazy oraz przetestowano proces wprowadzania danych testowych i wykonywania zaawansowanych zapytań SQL w dwóch środowiskach: \sphinxstylestrong{SQLite} oraz \sphinxstylestrong{PostgreSQL}.


\section{1.3. Wnioski końcowe z zajęć laboratoryjnych}
\label{\detokenize{rozdzial_1/index:wnioski-koncowe-z-zajec-laboratoryjnych}}
\sphinxAtStartPar
Przeprowadzone ćwiczenia i eksperymenty pozwoliły na wyciągnięcie następujących wniosków:
\begin{itemize}
\item {} 
\sphinxAtStartPar
\sphinxstylestrong{Projektowanie struktury:} Normalizacja do 3. postaci normalnej (3NF) jest kluczowa w systemach takich jak zarządzanie zleceniami, aby uniknąć redundancji danych i anomalii przy aktualizacji statusów zleceń lub magazynu.

\item {} 
\sphinxAtStartPar
\sphinxstylestrong{Proces importu danych:} Mechanizm \sphinxcode{\sphinxupquote{COPY}} dostępny w PostgreSQL znacząco przewyższa wydajnością wielokrotne instrukcje \sphinxcode{\sphinxupquote{INSERT}} w SQLite podczas wsadowego ładowania dużych zbiorów danych (np. tysięcy zużytych części).

\item {} 
\sphinxAtStartPar
\sphinxstylestrong{Różnice w dialektach SQL:} Choć oba silniki wspierają standard SQL, widoczne są różnice w zaawansowanych funkcjach wierszowych oraz mechanizmach auto\sphinxhyphen{}inkrementacji (\sphinxcode{\sphinxupquote{SERIAL}} w PostgreSQL vs zintegrowany mechanizm dla kluczy głównych w SQLite). PostgreSQL oferuje znacznie bogatszy system typów (np. precyzyjny typ \sphinxcode{\sphinxupquote{NUMERIC}} dla kosztów robocizny).

\end{itemize}

\sphinxstepscope


\chapter{Przegląd literatury}
\label{\detokenize{rozdzial_2/index:przeglad-literatury}}\label{\detokenize{rozdzial_2/index::doc}}
\sphinxstepscope

\sphinxAtStartPar
/Author: {[}Damian Dominiak Michał Bystrzak{]} / Date: 2026\sphinxhyphen{}05\sphinxhyphen{}14 / Version: 1.0 /


\chapter{Rozdział 3. Dokumentacja projektowa bazy danych}
\label{\detokenize{rozdzial_3/index:rozdzial-3-dokumentacja-projektowa-bazy-danych}}\label{\detokenize{rozdzial_3/index::doc}}

\section{3.1. Wybór zagadnienia i opis procesów}
\label{\detokenize{rozdzial_3/index:wybor-zagadnienia-i-opis-procesow}}
\sphinxAtStartPar
Tematem projektu jest relacyjna baza danych dla systemu zarządzania serwisem sprzętu IT. Aplikacja docelowa ma za zadanie obsługiwać pełen cykl życia zlecenia serwisowego \textendash{} od momentu przyjęcia sprzętu od klienta, poprzez diagnozę i zużycie części magazynowych, aż po ostateczne rozliczenie i wydanie naprawionego urządzenia.

\sphinxAtStartPar
\sphinxstylestrong{Główne procesy obsługiwane przez system:}
\begin{enumerate}
\sphinxsetlistlabels{\arabic}{enumi}{enumii}{}{.}%
\item {} 
\sphinxAtStartPar
\sphinxstylestrong{Ewidencja klientów i urządzeń:} Rejestracja danych kontaktowych klienta oraz powiązanie z nim konkretnych egzemplarzy sprzętu (identyfikowanych m.in. na podstawie numerów seryjnych).

\item {} 
\sphinxAtStartPar
\sphinxstylestrong{Zarządzanie zleceniami:} Tworzenie nowych zleceń naprawy dla konkretnych urządzeń, przypisywanie do nich techników prowadzących oraz śledzenie statusu prac.

\item {} 
\sphinxAtStartPar
\sphinxstylestrong{Gospodarka magazynowa części:} Rejestrowanie zużycia części wymiennych w ramach konkretnych zleceń, co pozwala na dokładne wyliczenie kosztów końcowych.

\end{enumerate}


\section{3.2. Prototypy danych i analiza typów}
\label{\detokenize{rozdzial_3/index:prototypy-danych-i-analiza-typow}}
\sphinxAtStartPar
W celu weryfikacji kompletności gromadzonych informacji przygotowano prototypowe pliki z jednym wierszem danych dla centralnej encji systemu, jaką jest zlecenie serwisowe.

\sphinxAtStartPar
\sphinxstylestrong{Prototyp w formacie JSON:}

\begin{sphinxVerbatim}[commandchars=\\\{\}]
\PYG{p}{\PYGZob{}}
\PYG{+w}{  }\PYG{n+nt}{\PYGZdq{}id\PYGZus{}zlecenia\PYGZdq{}}\PYG{p}{:}\PYG{+w}{ }\PYG{l+m+mi}{1024}\PYG{p}{,}
\PYG{+w}{  }\PYG{n+nt}{\PYGZdq{}data\PYGZus{}przyjecia\PYGZdq{}}\PYG{p}{:}\PYG{+w}{ }\PYG{l+s+s2}{\PYGZdq{}2024\PYGZhy{}05\PYGZhy{}14T12:00:00Z\PYGZdq{}}\PYG{p}{,}
\PYG{+w}{  }\PYG{n+nt}{\PYGZdq{}opis\PYGZus{}usterki\PYGZdq{}}\PYG{p}{:}\PYG{+w}{ }\PYG{l+s+s2}{\PYGZdq{}Uszkodzona matryca \PYGZhy{} pionowe pasy na ekranie\PYGZdq{}}\PYG{p}{,}
\PYG{+w}{  }\PYG{n+nt}{\PYGZdq{}status\PYGZdq{}}\PYG{p}{:}\PYG{+w}{ }\PYG{l+s+s2}{\PYGZdq{}W trakcie diagnozy\PYGZdq{}}\PYG{p}{,}
\PYG{+w}{  }\PYG{n+nt}{\PYGZdq{}koszt\PYGZus{}robocizny\PYGZdq{}}\PYG{p}{:}\PYG{+w}{ }\PYG{l+m+mf}{180.50}\PYG{p}{,}
\PYG{+w}{  }\PYG{n+nt}{\PYGZdq{}id\PYGZus{}klienta\PYGZdq{}}\PYG{p}{:}\PYG{+w}{ }\PYG{l+m+mi}{45}\PYG{p}{,}
\PYG{+w}{  }\PYG{n+nt}{\PYGZdq{}id\PYGZus{}technika\PYGZdq{}}\PYG{p}{:}\PYG{+w}{ }\PYG{l+m+mi}{12}
\PYG{p}{\PYGZcb{}}
\end{sphinxVerbatim}

\sphinxAtStartPar
\sphinxstylestrong{Prototyp w formacie CSV:}

\begin{sphinxVerbatim}[commandchars=\\\{\}]
id\PYGZus{}zlecenia,data\PYGZus{}przyjecia,opis\PYGZus{}usterki,status,koszt\PYGZus{}robocizny,id\PYGZus{}klienta,id\PYGZus{}technika
1024,2024\PYGZhy{}05\PYGZhy{}14T12:00:00Z,Uszkodzona matryca \PYGZhy{} pionowe pasy na ekranie,W trakcie diagnozy,180.50,45,12
\end{sphinxVerbatim}

\sphinxAtStartPar
\sphinxstylestrong{Zestawienie typów danych dla wybranych silników:}


\begin{savenotes}\sphinxattablestart
\sphinxthistablewithglobalstyle
\centering
\begin{tabular}[t]{\X{30}{100}\X{35}{100}\X{35}{100}}
\sphinxtoprule
\begin{varwidth}[t]{\sphinxcolwidth{1}{3}}
\sphinxstyletheadfamily \sphinxAtStartPar
Atrybut
\sphinxbeforeendvarwidth
\end{varwidth}%
&\begin{varwidth}[t]{\sphinxcolwidth{1}{3}}
\sphinxstyletheadfamily \sphinxAtStartPar
Typ danych (SQLite)
\sphinxbeforeendvarwidth
\end{varwidth}%
&\begin{varwidth}[t]{\sphinxcolwidth{1}{3}}
\sphinxstyletheadfamily \sphinxAtStartPar
Typ danych (PostgreSQL)
\sphinxbeforeendvarwidth
\end{varwidth}%
\\
\sphinxmidrule
\sphinxtableatstartofbodyhook\begin{varwidth}[t]{\sphinxcolwidth{1}{3}}
\sphinxAtStartPar
\sphinxcode{\sphinxupquote{id\_zlecenia}}
\sphinxbeforeendvarwidth
\end{varwidth}%
&\begin{varwidth}[t]{\sphinxcolwidth{1}{3}}
\sphinxAtStartPar
INTEGER PRIMARY KEY
\sphinxbeforeendvarwidth
\end{varwidth}%
&\begin{varwidth}[t]{\sphinxcolwidth{1}{3}}
\sphinxAtStartPar
SERIAL PRIMARY KEY
\sphinxbeforeendvarwidth
\end{varwidth}%
\\
\sphinxhline\begin{varwidth}[t]{\sphinxcolwidth{1}{3}}
\sphinxAtStartPar
\sphinxcode{\sphinxupquote{data\_przyjecia}}
\sphinxbeforeendvarwidth
\end{varwidth}%
&\begin{varwidth}[t]{\sphinxcolwidth{1}{3}}
\sphinxAtStartPar
TEXT
\sphinxbeforeendvarwidth
\end{varwidth}%
&\begin{varwidth}[t]{\sphinxcolwidth{1}{3}}
\sphinxAtStartPar
TIMESTAMP
\sphinxbeforeendvarwidth
\end{varwidth}%
\\
\sphinxhline\begin{varwidth}[t]{\sphinxcolwidth{1}{3}}
\sphinxAtStartPar
\sphinxcode{\sphinxupquote{opis\_usterki}}
\sphinxbeforeendvarwidth
\end{varwidth}%
&\begin{varwidth}[t]{\sphinxcolwidth{1}{3}}
\sphinxAtStartPar
TEXT
\sphinxbeforeendvarwidth
\end{varwidth}%
&\begin{varwidth}[t]{\sphinxcolwidth{1}{3}}
\sphinxAtStartPar
TEXT
\sphinxbeforeendvarwidth
\end{varwidth}%
\\
\sphinxhline\begin{varwidth}[t]{\sphinxcolwidth{1}{3}}
\sphinxAtStartPar
\sphinxcode{\sphinxupquote{status}}
\sphinxbeforeendvarwidth
\end{varwidth}%
&\begin{varwidth}[t]{\sphinxcolwidth{1}{3}}
\sphinxAtStartPar
TEXT
\sphinxbeforeendvarwidth
\end{varwidth}%
&\begin{varwidth}[t]{\sphinxcolwidth{1}{3}}
\sphinxAtStartPar
VARCHAR(50)
\sphinxbeforeendvarwidth
\end{varwidth}%
\\
\sphinxhline\begin{varwidth}[t]{\sphinxcolwidth{1}{3}}
\sphinxAtStartPar
\sphinxcode{\sphinxupquote{koszt\_robocizny}}
\sphinxbeforeendvarwidth
\end{varwidth}%
&\begin{varwidth}[t]{\sphinxcolwidth{1}{3}}
\sphinxAtStartPar
REAL
\sphinxbeforeendvarwidth
\end{varwidth}%
&\begin{varwidth}[t]{\sphinxcolwidth{1}{3}}
\sphinxAtStartPar
NUMERIC(10, 2)
\sphinxbeforeendvarwidth
\end{varwidth}%
\\
\sphinxhline\begin{varwidth}[t]{\sphinxcolwidth{1}{3}}
\sphinxAtStartPar
\sphinxcode{\sphinxupquote{id\_klienta}}
\sphinxbeforeendvarwidth
\end{varwidth}%
&\begin{varwidth}[t]{\sphinxcolwidth{1}{3}}
\sphinxAtStartPar
INTEGER
\sphinxbeforeendvarwidth
\end{varwidth}%
&\begin{varwidth}[t]{\sphinxcolwidth{1}{3}}
\sphinxAtStartPar
INT
\sphinxbeforeendvarwidth
\end{varwidth}%
\\
\sphinxbottomrule
\end{tabular}
\sphinxtableafterendhook\par
\sphinxattableend\end{savenotes}


\section{3.3. Model koncepcyjny bazy danych}
\label{\detokenize{rozdzial_3/index:model-koncepcyjny-bazy-danych}}
\sphinxAtStartPar
Na etapie modelowania koncepcyjnego zdefiniowano główne wymagania informacyjne systemu. Zidentyfikowano encje, ich atrybuty oraz relacje zachodzące w analizowanym środowisku. Do wizualizacji wybrano notację Crow’s Foot (kurzych łapek) jako współczesny standard branżowy, który zapewnia wyższą czytelność i płynniejsze przejście do modelu logicznego niż tradycyjna notacja Chena.

\sphinxAtStartPar
\sphinxstylestrong{Zidentyfikowane encje mocne i ich klucze:}
\begin{itemize}
\item {} 
\sphinxAtStartPar
\sphinxstylestrong{KLIENT:} \sphinxcode{\sphinxupquote{ID\_Klienta}} (PK), Imię, Nazwisko, Telefon, Email.

\item {} 
\sphinxAtStartPar
\sphinxstylestrong{PRACOWNIK:} \sphinxcode{\sphinxupquote{ID\_Pracownika}} (PK), Imię, Nazwisko, Stanowisko.

\item {} 
\sphinxAtStartPar
\sphinxstylestrong{ZLECENIE:} \sphinxcode{\sphinxupquote{ID\_Zlecenia}} (PK), Data\_przyjecia, Opis\_usterki, Status, Koszt\_robocizny.

\item {} 
\sphinxAtStartPar
\sphinxstylestrong{CZĘŚĆ:} \sphinxcode{\sphinxupquote{ID\_Czesci}} (PK), Nazwa, Producent, Stan\_magazynowy, Cena.

\end{itemize}

\sphinxAtStartPar
\sphinxstylestrong{Zidentyfikowana encja słaba:}
\begin{itemize}
\item {} 
\sphinxAtStartPar
\sphinxstylestrong{URZĄDZENIE:} Posiada klucz częściowy \sphinxcode{\sphinxupquote{Numer\_kolejny}} oraz atrybuty Marka, Model, Numer\_seryjny. Jest to encja słaba, ponieważ jej egzystencja w systemie jest całkowicie zależna od istnienia powiązanej encji mocnej (KLIENT).

\end{itemize}

\sphinxAtStartPar
\sphinxstylestrong{Związki i krotności:}

\sphinxAtStartPar
W modelu poprawnie zamodelowano ścieżkę powiązań \sphinxcode{\sphinxupquote{KLIENT}} \sphinxhyphen{}\textgreater{} \sphinxcode{\sphinxupquote{URZĄDZENIE}} \sphinxhyphen{}\textgreater{} \sphinxcode{\sphinxupquote{ZLECENIE}}. Uniknięto w ten sposób związku niepoprawnego, zwanego pułapką wiatraka (fan trap), który wystąpiłby, gdyby encje \sphinxcode{\sphinxupquote{URZĄDZENIE}} oraz \sphinxcode{\sphinxupquote{ZLECENIE}} były połączone z encją \sphinxcode{\sphinxupquote{KLIENT}} w sposób niezależny.
W modelu zidentyfikowano również związek wiele\sphinxhyphen{}do\sphinxhyphen{}wielu (M:N) pomiędzy encjami \sphinxcode{\sphinxupquote{ZLECENIE}} oraz \sphinxcode{\sphinxupquote{CZĘŚĆ}}.

\begin{figure}[htbp]
\centering
\capstart

\noindent\sphinxincludegraphics{{diagram_koncepcyjny}.jpeg}
\caption{Rys. 1. Diagram koncepcyjny bazy danych serwisu IT.}\label{\detokenize{rozdzial_3/index:id1}}\end{figure}


\section{3.4. Model logiczny i normalizacja}
\label{\detokenize{rozdzial_3/index:model-logiczny-i-normalizacja}}
\sphinxAtStartPar
Model logiczny został opracowany na bazie modelu koncepcyjnego poprzez uszczegółowienie struktury i dostosowanie jej do paradygmatu relacyjnego.

\sphinxAtStartPar
\sphinxstylestrong{Eliminacja relacji M:N}

\sphinxAtStartPar
Występująca w modelu koncepcyjnym relacja wiele\sphinxhyphen{}do\sphinxhyphen{}wielu między zleceniami a częściami została rozwiązana poprzez wprowadzenie tabeli asocjacyjnej (pośredniczącej) o nazwie \sphinxcode{\sphinxupquote{ZUZYTE\_CZESCI}}. Tabela ta posiada złożony klucz główny (\sphinxcode{\sphinxupquote{ID\_Zlecenia}}, \sphinxcode{\sphinxupquote{ID\_Czesci}}) oraz dodatkowy atrybut niekluczowy \sphinxcode{\sphinxupquote{Ilosc}}. Wprowadzono również klucze sztuczne i obce.

\sphinxAtStartPar
\sphinxstylestrong{Proces normalizacji:}

\sphinxAtStartPar
Zaprojektowana struktura została poddana procesowi normalizacji, osiągając \sphinxstylestrong{3. postać normalną (3NF)}, co stanowi optymalny kompromis między wydajnością systemu a eliminacją redundancji danych.
\begin{itemize}
\item {} 
\sphinxAtStartPar
\sphinxstylestrong{1NF:} Tabela spełnia pierwszą postać normalną, ponieważ wszystkie atrybuty są atomowe, a dane nadmiarowe (np. wiele użytych części w jednym zleceniu) zostały przeniesione do nowej tabeli.

\item {} 
\sphinxAtStartPar
\sphinxstylestrong{2NF:} Warunek drugiej postaci normalnej jest spełniony. Każdy atrybut niekluczowy jest w pełni funkcyjnie zależny od całego klucza głównego (istotne zwłaszcza w tabeli \sphinxcode{\sphinxupquote{ZUZYTE\_CZESCI}}, gdzie wartość \sphinxcode{\sphinxupquote{Ilosc}} zależy wprost od pary zlecenie\sphinxhyphen{}część).

\item {} 
\sphinxAtStartPar
\sphinxstylestrong{3NF:} Model osiąga trzecią postać normalną, gdyż nie występują w nim żadne zależności przechodnie (tranzytywne) między atrybutami niekluczowymi.

\end{itemize}

\begin{figure}[htbp]
\centering
\capstart

\noindent\sphinxincludegraphics{{diagram_logiczny}.jpeg}
\caption{Rys. 2. Diagram logiczny bazy danych serwisu IT w 3NF.}\label{\detokenize{rozdzial_3/index:id2}}\end{figure}


\section{3.5. Model fizyczny}
\label{\detokenize{rozdzial_3/index:model-fizyczny}}
\sphinxAtStartPar
Na podstawie schematu logicznego przygotowano fizyczną implementację struktury danych z uwzględnieniem specyfiki dwóch różnych systemów zarządzania relacyjnymi bazami danych: SQLite oraz PostgreSQL.


\subsection{3.5.1. Implementacja w środowisku SQLite}
\label{\detokenize{rozdzial_3/index:implementacja-w-srodowisku-sqlite}}
\sphinxAtStartPar
Model wykorzystuje proste typowanie dostępne w silniku SQLite.

\begin{sphinxVerbatim}[commandchars=\\\{\}]
\PYG{k}{CREATE}\PYG{+w}{ }\PYG{k}{TABLE}\PYG{+w}{ }\PYG{n}{KLIENT}\PYG{+w}{ }\PYG{p}{(}
\PYG{+w}{    }\PYG{n}{ID\PYGZus{}Klienta}\PYG{+w}{ }\PYG{n+nb}{INTEGER}\PYG{+w}{ }\PYG{k}{PRIMARY}\PYG{+w}{ }\PYG{k}{KEY}\PYG{p}{,}
\PYG{+w}{    }\PYG{n}{Imie}\PYG{+w}{ }\PYG{n+nb}{TEXT}\PYG{+w}{ }\PYG{k}{NOT}\PYG{+w}{ }\PYG{k}{NULL}\PYG{p}{,}
\PYG{+w}{    }\PYG{n}{Nazwisko}\PYG{+w}{ }\PYG{n+nb}{TEXT}\PYG{+w}{ }\PYG{k}{NOT}\PYG{+w}{ }\PYG{k}{NULL}\PYG{p}{,}
\PYG{+w}{    }\PYG{n}{Telefon}\PYG{+w}{ }\PYG{n+nb}{TEXT}\PYG{p}{,}
\PYG{+w}{    }\PYG{n}{Email}\PYG{+w}{ }\PYG{n+nb}{TEXT}
\PYG{p}{)}\PYG{p}{;}

\PYG{k}{CREATE}\PYG{+w}{ }\PYG{k}{TABLE}\PYG{+w}{ }\PYG{n}{URZADZENIE}\PYG{+w}{ }\PYG{p}{(}
\PYG{+w}{    }\PYG{n}{ID\PYGZus{}Urzadzenia}\PYG{+w}{ }\PYG{n+nb}{INTEGER}\PYG{+w}{ }\PYG{k}{PRIMARY}\PYG{+w}{ }\PYG{k}{KEY}\PYG{p}{,}
\PYG{+w}{    }\PYG{n}{ID\PYGZus{}Klienta}\PYG{+w}{ }\PYG{n+nb}{INTEGER}\PYG{+w}{ }\PYG{k}{NOT}\PYG{+w}{ }\PYG{k}{NULL}\PYG{p}{,}
\PYG{+w}{    }\PYG{n}{Marka}\PYG{+w}{ }\PYG{n+nb}{TEXT}\PYG{+w}{ }\PYG{k}{NOT}\PYG{+w}{ }\PYG{k}{NULL}\PYG{p}{,}
\PYG{+w}{    }\PYG{n}{Model}\PYG{+w}{ }\PYG{n+nb}{TEXT}\PYG{+w}{ }\PYG{k}{NOT}\PYG{+w}{ }\PYG{k}{NULL}\PYG{p}{,}
\PYG{+w}{    }\PYG{n}{Numer\PYGZus{}seryjny}\PYG{+w}{ }\PYG{n+nb}{TEXT}\PYG{p}{,}
\PYG{+w}{    }\PYG{k}{FOREIGN}\PYG{+w}{ }\PYG{k}{KEY}\PYG{+w}{ }\PYG{p}{(}\PYG{n}{ID\PYGZus{}Klienta}\PYG{p}{)}\PYG{+w}{ }\PYG{k}{REFERENCES}\PYG{+w}{ }\PYG{n}{KLIENT}\PYG{p}{(}\PYG{n}{ID\PYGZus{}Klienta}\PYG{p}{)}
\PYG{p}{)}\PYG{p}{;}

\PYG{k}{CREATE}\PYG{+w}{ }\PYG{k}{TABLE}\PYG{+w}{ }\PYG{n}{PRACOWNIK}\PYG{+w}{ }\PYG{p}{(}
\PYG{+w}{    }\PYG{n}{ID\PYGZus{}Pracownika}\PYG{+w}{ }\PYG{n+nb}{INTEGER}\PYG{+w}{ }\PYG{k}{PRIMARY}\PYG{+w}{ }\PYG{k}{KEY}\PYG{p}{,}
\PYG{+w}{    }\PYG{n}{Imie}\PYG{+w}{ }\PYG{n+nb}{TEXT}\PYG{+w}{ }\PYG{k}{NOT}\PYG{+w}{ }\PYG{k}{NULL}\PYG{p}{,}
\PYG{+w}{    }\PYG{n}{Nazwisko}\PYG{+w}{ }\PYG{n+nb}{TEXT}\PYG{+w}{ }\PYG{k}{NOT}\PYG{+w}{ }\PYG{k}{NULL}\PYG{p}{,}
\PYG{+w}{    }\PYG{n}{Stanowisko}\PYG{+w}{ }\PYG{n+nb}{TEXT}\PYG{+w}{ }\PYG{k}{NOT}\PYG{+w}{ }\PYG{k}{NULL}
\PYG{p}{)}\PYG{p}{;}

\PYG{k}{CREATE}\PYG{+w}{ }\PYG{k}{TABLE}\PYG{+w}{ }\PYG{n}{ZLECENIE}\PYG{+w}{ }\PYG{p}{(}
\PYG{+w}{    }\PYG{n}{ID\PYGZus{}Zlecenia}\PYG{+w}{ }\PYG{n+nb}{INTEGER}\PYG{+w}{ }\PYG{k}{PRIMARY}\PYG{+w}{ }\PYG{k}{KEY}\PYG{p}{,}
\PYG{+w}{    }\PYG{n}{ID\PYGZus{}Urzadzenia}\PYG{+w}{ }\PYG{n+nb}{INTEGER}\PYG{+w}{ }\PYG{k}{NOT}\PYG{+w}{ }\PYG{k}{NULL}\PYG{p}{,}
\PYG{+w}{    }\PYG{n}{ID\PYGZus{}Pracownika}\PYG{+w}{ }\PYG{n+nb}{INTEGER}\PYG{+w}{ }\PYG{k}{NOT}\PYG{+w}{ }\PYG{k}{NULL}\PYG{p}{,}
\PYG{+w}{    }\PYG{n}{Data\PYGZus{}przyjecia}\PYG{+w}{ }\PYG{n+nb}{TEXT}\PYG{+w}{ }\PYG{k}{NOT}\PYG{+w}{ }\PYG{k}{NULL}\PYG{p}{,}
\PYG{+w}{    }\PYG{n}{Opis\PYGZus{}usterki}\PYG{+w}{ }\PYG{n+nb}{TEXT}\PYG{+w}{ }\PYG{k}{NOT}\PYG{+w}{ }\PYG{k}{NULL}\PYG{p}{,}
\PYG{+w}{    }\PYG{n}{Status}\PYG{+w}{ }\PYG{n+nb}{TEXT}\PYG{+w}{ }\PYG{k}{NOT}\PYG{+w}{ }\PYG{k}{NULL}\PYG{p}{,}
\PYG{+w}{    }\PYG{n}{Koszt\PYGZus{}robocizny}\PYG{+w}{ }\PYG{n+nb}{REAL}\PYG{p}{,}
\PYG{+w}{    }\PYG{k}{FOREIGN}\PYG{+w}{ }\PYG{k}{KEY}\PYG{+w}{ }\PYG{p}{(}\PYG{n}{ID\PYGZus{}Urzadzenia}\PYG{p}{)}\PYG{+w}{ }\PYG{k}{REFERENCES}\PYG{+w}{ }\PYG{n}{URZADZENIE}\PYG{p}{(}\PYG{n}{ID\PYGZus{}Urzadzenia}\PYG{p}{)}\PYG{p}{,}
\PYG{+w}{    }\PYG{k}{FOREIGN}\PYG{+w}{ }\PYG{k}{KEY}\PYG{+w}{ }\PYG{p}{(}\PYG{n}{ID\PYGZus{}Pracownika}\PYG{p}{)}\PYG{+w}{ }\PYG{k}{REFERENCES}\PYG{+w}{ }\PYG{n}{PRACOWNIK}\PYG{p}{(}\PYG{n}{ID\PYGZus{}Pracownika}\PYG{p}{)}
\PYG{p}{)}\PYG{p}{;}

\PYG{k}{CREATE}\PYG{+w}{ }\PYG{k}{TABLE}\PYG{+w}{ }\PYG{n}{CZESC}\PYG{+w}{ }\PYG{p}{(}
\PYG{+w}{    }\PYG{n}{ID\PYGZus{}Czesci}\PYG{+w}{ }\PYG{n+nb}{INTEGER}\PYG{+w}{ }\PYG{k}{PRIMARY}\PYG{+w}{ }\PYG{k}{KEY}\PYG{p}{,}
\PYG{+w}{    }\PYG{n}{Nazwa}\PYG{+w}{ }\PYG{n+nb}{TEXT}\PYG{+w}{ }\PYG{k}{NOT}\PYG{+w}{ }\PYG{k}{NULL}\PYG{p}{,}
\PYG{+w}{    }\PYG{n}{Producent}\PYG{+w}{ }\PYG{n+nb}{TEXT}\PYG{p}{,}
\PYG{+w}{    }\PYG{n}{Stan\PYGZus{}magazynowy}\PYG{+w}{ }\PYG{n+nb}{INTEGER}\PYG{+w}{ }\PYG{k}{DEFAULT}\PYG{+w}{ }\PYG{l+m+mi}{0}\PYG{p}{,}
\PYG{+w}{    }\PYG{n}{Cena}\PYG{+w}{ }\PYG{n+nb}{REAL}\PYG{+w}{ }\PYG{k}{NOT}\PYG{+w}{ }\PYG{k}{NULL}
\PYG{p}{)}\PYG{p}{;}

\PYG{k}{CREATE}\PYG{+w}{ }\PYG{k}{TABLE}\PYG{+w}{ }\PYG{n}{ZUZYTE\PYGZus{}CZESCI}\PYG{+w}{ }\PYG{p}{(}
\PYG{+w}{    }\PYG{n}{ID\PYGZus{}Zlecenia}\PYG{+w}{ }\PYG{n+nb}{INTEGER}\PYG{p}{,}
\PYG{+w}{    }\PYG{n}{ID\PYGZus{}Czesci}\PYG{+w}{ }\PYG{n+nb}{INTEGER}\PYG{p}{,}
\PYG{+w}{    }\PYG{n}{Ilosc}\PYG{+w}{ }\PYG{n+nb}{INTEGER}\PYG{+w}{ }\PYG{k}{NOT}\PYG{+w}{ }\PYG{k}{NULL}\PYG{+w}{ }\PYG{k}{DEFAULT}\PYG{+w}{ }\PYG{l+m+mi}{1}\PYG{p}{,}
\PYG{+w}{    }\PYG{k}{PRIMARY}\PYG{+w}{ }\PYG{k}{KEY}\PYG{+w}{ }\PYG{p}{(}\PYG{n}{ID\PYGZus{}Zlecenia}\PYG{p}{,}\PYG{+w}{ }\PYG{n}{ID\PYGZus{}Czesci}\PYG{p}{)}\PYG{p}{,}
\PYG{+w}{    }\PYG{k}{FOREIGN}\PYG{+w}{ }\PYG{k}{KEY}\PYG{+w}{ }\PYG{p}{(}\PYG{n}{ID\PYGZus{}Zlecenia}\PYG{p}{)}\PYG{+w}{ }\PYG{k}{REFERENCES}\PYG{+w}{ }\PYG{n}{ZLECENIE}\PYG{p}{(}\PYG{n}{ID\PYGZus{}Zlecenia}\PYG{p}{)}\PYG{p}{,}
\PYG{+w}{    }\PYG{k}{FOREIGN}\PYG{+w}{ }\PYG{k}{KEY}\PYG{+w}{ }\PYG{p}{(}\PYG{n}{ID\PYGZus{}Czesci}\PYG{p}{)}\PYG{+w}{ }\PYG{k}{REFERENCES}\PYG{+w}{ }\PYG{n}{CZESC}\PYG{p}{(}\PYG{n}{ID\PYGZus{}Czesci}\PYG{p}{)}
\PYG{p}{)}\PYG{p}{;}
\end{sphinxVerbatim}


\subsection{3.5.2. Implementacja w środowisku PostgreSQL}
\label{\detokenize{rozdzial_3/index:implementacja-w-srodowisku-postgresql}}
\sphinxAtStartPar
Model dostosowany do zaawansowanych możliwości silnika PostgreSQL, wykorzystujący automatyczne generowanie kluczy głównych (typ \sphinxcode{\sphinxupquote{SERIAL}}), precyzyjne typowanie numeryczne (\sphinxcode{\sphinxupquote{NUMERIC}}) oraz rygorystyczne więzy integralności (np. \sphinxcode{\sphinxupquote{CHECK}}).

\begin{sphinxVerbatim}[commandchars=\\\{\}]
\PYG{k}{CREATE}\PYG{+w}{ }\PYG{k}{TABLE}\PYG{+w}{ }\PYG{n}{KLIENT}\PYG{+w}{ }\PYG{p}{(}
\PYG{+w}{    }\PYG{n}{ID\PYGZus{}Klienta}\PYG{+w}{ }\PYG{n+nb}{SERIAL}\PYG{+w}{ }\PYG{k}{PRIMARY}\PYG{+w}{ }\PYG{k}{KEY}\PYG{p}{,}
\PYG{+w}{    }\PYG{n}{Imie}\PYG{+w}{ }\PYG{n+nb}{VARCHAR}\PYG{p}{(}\PYG{l+m+mi}{50}\PYG{p}{)}\PYG{+w}{ }\PYG{k}{NOT}\PYG{+w}{ }\PYG{k}{NULL}\PYG{p}{,}
\PYG{+w}{    }\PYG{n}{Nazwisko}\PYG{+w}{ }\PYG{n+nb}{VARCHAR}\PYG{p}{(}\PYG{l+m+mi}{50}\PYG{p}{)}\PYG{+w}{ }\PYG{k}{NOT}\PYG{+w}{ }\PYG{k}{NULL}\PYG{p}{,}
\PYG{+w}{    }\PYG{n}{Telefon}\PYG{+w}{ }\PYG{n+nb}{VARCHAR}\PYG{p}{(}\PYG{l+m+mi}{15}\PYG{p}{)}\PYG{p}{,}
\PYG{+w}{    }\PYG{n}{Email}\PYG{+w}{ }\PYG{n+nb}{VARCHAR}\PYG{p}{(}\PYG{l+m+mi}{100}\PYG{p}{)}
\PYG{p}{)}\PYG{p}{;}

\PYG{k}{CREATE}\PYG{+w}{ }\PYG{k}{TABLE}\PYG{+w}{ }\PYG{n}{URZADZENIE}\PYG{+w}{ }\PYG{p}{(}
\PYG{+w}{    }\PYG{n}{ID\PYGZus{}Urzadzenia}\PYG{+w}{ }\PYG{n+nb}{SERIAL}\PYG{+w}{ }\PYG{k}{PRIMARY}\PYG{+w}{ }\PYG{k}{KEY}\PYG{p}{,}
\PYG{+w}{    }\PYG{n}{ID\PYGZus{}Klienta}\PYG{+w}{ }\PYG{n+nb}{INTEGER}\PYG{+w}{ }\PYG{k}{NOT}\PYG{+w}{ }\PYG{k}{NULL}\PYG{+w}{ }\PYG{k}{REFERENCES}\PYG{+w}{ }\PYG{n}{KLIENT}\PYG{p}{(}\PYG{n}{ID\PYGZus{}Klienta}\PYG{p}{)}\PYG{p}{,}
\PYG{+w}{    }\PYG{n}{Marka}\PYG{+w}{ }\PYG{n+nb}{VARCHAR}\PYG{p}{(}\PYG{l+m+mi}{50}\PYG{p}{)}\PYG{+w}{ }\PYG{k}{NOT}\PYG{+w}{ }\PYG{k}{NULL}\PYG{p}{,}
\PYG{+w}{    }\PYG{n}{Model}\PYG{+w}{ }\PYG{n+nb}{VARCHAR}\PYG{p}{(}\PYG{l+m+mi}{50}\PYG{p}{)}\PYG{+w}{ }\PYG{k}{NOT}\PYG{+w}{ }\PYG{k}{NULL}\PYG{p}{,}
\PYG{+w}{    }\PYG{n}{Numer\PYGZus{}seryjny}\PYG{+w}{ }\PYG{n+nb}{VARCHAR}\PYG{p}{(}\PYG{l+m+mi}{100}\PYG{p}{)}\PYG{+w}{ }\PYG{k}{UNIQUE}
\PYG{p}{)}\PYG{p}{;}

\PYG{k}{CREATE}\PYG{+w}{ }\PYG{k}{TABLE}\PYG{+w}{ }\PYG{n}{PRACOWNIK}\PYG{+w}{ }\PYG{p}{(}
\PYG{+w}{    }\PYG{n}{ID\PYGZus{}Pracownika}\PYG{+w}{ }\PYG{n+nb}{SERIAL}\PYG{+w}{ }\PYG{k}{PRIMARY}\PYG{+w}{ }\PYG{k}{KEY}\PYG{p}{,}
\PYG{+w}{    }\PYG{n}{Imie}\PYG{+w}{ }\PYG{n+nb}{VARCHAR}\PYG{p}{(}\PYG{l+m+mi}{50}\PYG{p}{)}\PYG{+w}{ }\PYG{k}{NOT}\PYG{+w}{ }\PYG{k}{NULL}\PYG{p}{,}
\PYG{+w}{    }\PYG{n}{Nazwisko}\PYG{+w}{ }\PYG{n+nb}{VARCHAR}\PYG{p}{(}\PYG{l+m+mi}{50}\PYG{p}{)}\PYG{+w}{ }\PYG{k}{NOT}\PYG{+w}{ }\PYG{k}{NULL}\PYG{p}{,}
\PYG{+w}{    }\PYG{n}{Stanowisko}\PYG{+w}{ }\PYG{n+nb}{VARCHAR}\PYG{p}{(}\PYG{l+m+mi}{50}\PYG{p}{)}\PYG{+w}{ }\PYG{k}{NOT}\PYG{+w}{ }\PYG{k}{NULL}
\PYG{p}{)}\PYG{p}{;}

\PYG{k}{CREATE}\PYG{+w}{ }\PYG{k}{TABLE}\PYG{+w}{ }\PYG{n}{ZLECENIE}\PYG{+w}{ }\PYG{p}{(}
\PYG{+w}{    }\PYG{n}{ID\PYGZus{}Zlecenia}\PYG{+w}{ }\PYG{n+nb}{SERIAL}\PYG{+w}{ }\PYG{k}{PRIMARY}\PYG{+w}{ }\PYG{k}{KEY}\PYG{p}{,}
\PYG{+w}{    }\PYG{n}{ID\PYGZus{}Urzadzenia}\PYG{+w}{ }\PYG{n+nb}{INTEGER}\PYG{+w}{ }\PYG{k}{NOT}\PYG{+w}{ }\PYG{k}{NULL}\PYG{+w}{ }\PYG{k}{REFERENCES}\PYG{+w}{ }\PYG{n}{URZADZENIE}\PYG{p}{(}\PYG{n}{ID\PYGZus{}Urzadzenia}\PYG{p}{)}\PYG{p}{,}
\PYG{+w}{    }\PYG{n}{ID\PYGZus{}Pracownika}\PYG{+w}{ }\PYG{n+nb}{INTEGER}\PYG{+w}{ }\PYG{k}{NOT}\PYG{+w}{ }\PYG{k}{NULL}\PYG{+w}{ }\PYG{k}{REFERENCES}\PYG{+w}{ }\PYG{n}{PRACOWNIK}\PYG{p}{(}\PYG{n}{ID\PYGZus{}Pracownika}\PYG{p}{)}\PYG{p}{,}
\PYG{+w}{    }\PYG{n}{Data\PYGZus{}przyjecia}\PYG{+w}{ }\PYG{k}{TIMESTAMP}\PYG{+w}{ }\PYG{k}{NOT}\PYG{+w}{ }\PYG{k}{NULL}\PYG{+w}{ }\PYG{k}{DEFAULT}\PYG{+w}{ }\PYG{k}{CURRENT\PYGZus{}TIMESTAMP}\PYG{p}{,}
\PYG{+w}{    }\PYG{n}{Opis\PYGZus{}usterki}\PYG{+w}{ }\PYG{n+nb}{TEXT}\PYG{+w}{ }\PYG{k}{NOT}\PYG{+w}{ }\PYG{k}{NULL}\PYG{p}{,}
\PYG{+w}{    }\PYG{n}{Status}\PYG{+w}{ }\PYG{n+nb}{VARCHAR}\PYG{p}{(}\PYG{l+m+mi}{30}\PYG{p}{)}\PYG{+w}{ }\PYG{k}{NOT}\PYG{+w}{ }\PYG{k}{NULL}\PYG{p}{,}
\PYG{+w}{    }\PYG{n}{Koszt\PYGZus{}robocizny}\PYG{+w}{ }\PYG{n+nb}{NUMERIC}\PYG{p}{(}\PYG{l+m+mi}{10}\PYG{p}{,}\PYG{+w}{ }\PYG{l+m+mi}{2}\PYG{p}{)}
\PYG{p}{)}\PYG{p}{;}

\PYG{k}{CREATE}\PYG{+w}{ }\PYG{k}{TABLE}\PYG{+w}{ }\PYG{n}{CZESC}\PYG{+w}{ }\PYG{p}{(}
\PYG{+w}{    }\PYG{n}{ID\PYGZus{}Czesci}\PYG{+w}{ }\PYG{n+nb}{SERIAL}\PYG{+w}{ }\PYG{k}{PRIMARY}\PYG{+w}{ }\PYG{k}{KEY}\PYG{p}{,}
\PYG{+w}{    }\PYG{n}{Nazwa}\PYG{+w}{ }\PYG{n+nb}{VARCHAR}\PYG{p}{(}\PYG{l+m+mi}{100}\PYG{p}{)}\PYG{+w}{ }\PYG{k}{NOT}\PYG{+w}{ }\PYG{k}{NULL}\PYG{p}{,}
\PYG{+w}{    }\PYG{n}{Producent}\PYG{+w}{ }\PYG{n+nb}{VARCHAR}\PYG{p}{(}\PYG{l+m+mi}{50}\PYG{p}{)}\PYG{p}{,}
\PYG{+w}{    }\PYG{n}{Stan\PYGZus{}magazynowy}\PYG{+w}{ }\PYG{n+nb}{INTEGER}\PYG{+w}{ }\PYG{k}{DEFAULT}\PYG{+w}{ }\PYG{l+m+mi}{0}\PYG{+w}{ }\PYG{k}{CHECK}\PYG{+w}{ }\PYG{p}{(}\PYG{n}{Stan\PYGZus{}magazynowy}\PYG{+w}{ }\PYG{o}{\PYGZgt{}}\PYG{o}{=}\PYG{+w}{ }\PYG{l+m+mi}{0}\PYG{p}{)}\PYG{p}{,}
\PYG{+w}{    }\PYG{n}{Cena}\PYG{+w}{ }\PYG{n+nb}{NUMERIC}\PYG{p}{(}\PYG{l+m+mi}{10}\PYG{p}{,}\PYG{+w}{ }\PYG{l+m+mi}{2}\PYG{p}{)}\PYG{+w}{ }\PYG{k}{NOT}\PYG{+w}{ }\PYG{k}{NULL}\PYG{+w}{ }\PYG{k}{CHECK}\PYG{+w}{ }\PYG{p}{(}\PYG{n}{Cena}\PYG{+w}{ }\PYG{o}{\PYGZgt{}}\PYG{o}{=}\PYG{+w}{ }\PYG{l+m+mi}{0}\PYG{p}{)}
\PYG{p}{)}\PYG{p}{;}

\PYG{k}{CREATE}\PYG{+w}{ }\PYG{k}{TABLE}\PYG{+w}{ }\PYG{n}{ZUZYTE\PYGZus{}CZESCI}\PYG{+w}{ }\PYG{p}{(}
\PYG{+w}{    }\PYG{n}{ID\PYGZus{}Zlecenia}\PYG{+w}{ }\PYG{n+nb}{INTEGER}\PYG{+w}{ }\PYG{k}{REFERENCES}\PYG{+w}{ }\PYG{n}{ZLECENIE}\PYG{p}{(}\PYG{n}{ID\PYGZus{}Zlecenia}\PYG{p}{)}\PYG{p}{,}
\PYG{+w}{    }\PYG{n}{ID\PYGZus{}Czesci}\PYG{+w}{ }\PYG{n+nb}{INTEGER}\PYG{+w}{ }\PYG{k}{REFERENCES}\PYG{+w}{ }\PYG{n}{CZESC}\PYG{p}{(}\PYG{n}{ID\PYGZus{}Czesci}\PYG{p}{)}\PYG{p}{,}
\PYG{+w}{    }\PYG{n}{Ilosc}\PYG{+w}{ }\PYG{n+nb}{INTEGER}\PYG{+w}{ }\PYG{k}{NOT}\PYG{+w}{ }\PYG{k}{NULL}\PYG{+w}{ }\PYG{k}{DEFAULT}\PYG{+w}{ }\PYG{l+m+mi}{1}\PYG{+w}{ }\PYG{k}{CHECK}\PYG{+w}{ }\PYG{p}{(}\PYG{n}{Ilosc}\PYG{+w}{ }\PYG{o}{\PYGZgt{}}\PYG{+w}{ }\PYG{l+m+mi}{0}\PYG{p}{)}\PYG{p}{,}
\PYG{+w}{    }\PYG{k}{PRIMARY}\PYG{+w}{ }\PYG{k}{KEY}\PYG{+w}{ }\PYG{p}{(}\PYG{n}{ID\PYGZus{}Zlecenia}\PYG{p}{,}\PYG{+w}{ }\PYG{n}{ID\PYGZus{}Czesci}\PYG{p}{)}
\PYG{p}{)}\PYG{p}{;}
\end{sphinxVerbatim}

\sphinxstepscope

\sphinxAtStartPar
/Author: {[}Damian Dominiak, Michał Bystrzak{]} / Date: 2026\sphinxhyphen{}05\sphinxhyphen{}15 / Version: 1.0 /


\chapter{Rozdział 4. Definiowanie bazy i wprowadzanie danych}
\label{\detokenize{rozdzial_4/index:rozdzial-4-definiowanie-bazy-i-wprowadzanie-danych}}\label{\detokenize{rozdzial_4/index::doc}}
\sphinxAtStartPar
W tym rozdziale opisano mechanizmy wykorzystane do wsadowego wprowadzania danych testowych do baz SQLite oraz PostgreSQL dla systemu serwisu IT.


\section{4.1. Generowanie danych testowych}
\label{\detokenize{rozdzial_4/index:generowanie-danych-testowych}}
\sphinxAtStartPar
W celu symulacji rzeczywistego środowiska pracy serwisu sprzętu IT, przygotowano skrypt w języku Python, który automatycznie generuje losowe dane testowe (m.in. dane klientów, urządzeń, pracowników oraz zleceń i zużytych części) i eksportuje je do plików w formacie CSV.


\section{4.2. Import danych do PostgreSQL}
\label{\detokenize{rozdzial_4/index:import-danych-do-postgresql}}
\sphinxAtStartPar
Dla silnika PostgreSQL zdecydowano się na wykorzystanie natywnego mechanizmu \sphinxcode{\sphinxupquote{COPY}}. Pozwala on na najszybszy bulk\sphinxhyphen{}insert danych z plików CSV.


\section{4.3. Import danych do SQLite}
\label{\detokenize{rozdzial_4/index:import-danych-do-sqlite}}
\sphinxAtStartPar
W przypadku bazy SQLite zaimplementowano skrypt w języku Python z wykorzystaniem wbudowanej biblioteki \sphinxcode{\sphinxupquote{sqlite3}}, który wczytuje dane z plików CSV i wykonuje instrukcje \sphinxcode{\sphinxupquote{INSERT}} wewnątrz transakcji, co znacząco optymalizuje czas wstawiania setek rekordów.

\sphinxstepscope

\sphinxAtStartPar
/Author: {[}Damian Dominiak, Michał Bystrzak{]} / Date: 2026\sphinxhyphen{}05\sphinxhyphen{}15 / Version: 1.0 /


\chapter{Rozdział 5. Zapytania do bazy danych}
\label{\detokenize{rozdzial_5/index:rozdzial-5-zapytania-do-bazy-danych}}\label{\detokenize{rozdzial_5/index::doc}}
\sphinxAtStartPar
Poniższa dokumentacja zawiera opisy funkcji odpowiedzialnych za wykonywanie zapytań do bazy danych (PostgreSQL oraz SQLite). Dokumentacja ta została wygenerowana automatycznie z komentarzy (docstrings) zawartych w module Pythona.



\renewcommand{\indexname}{Indeks}
\printindex
\end{document}